% index.tex -- index of Module 7 multiple-choice questions in the 2019-2025 NSW HSC Physics papers.
% Descriptions are paraphrases, not question text.
\documentclass[12pt]{article}
\usepackage{quiz}
\usepackage{longtable}
\hypersetup{pdftitle={HSC Physics Module 7 Multiple-Choice Questions 2019-2025 - Index}}
\begin{document}
\quiztitle{Year 12 Physics}{HSC Module 7 Multiple-Choice Questions 2019–2025: Index}
Every Section I question from the 2019--2025 NSW HSC Physics papers that tests Module 7
(The Nature of Light), with the answer from the NESA marking guidelines. Topics: T1 electromagnetic spectrum; T2 light: wave model; T3 light: quantum model; T4 light and special relativity. A second module after
the topic (for example T3/M8) means the question also uses ideas from that module. The descriptions are
short summaries, not the question text.

{\small\renewcommand\arraystretch{1.3}
\begin{longtable}{|c|c|c|p{108mm}|c|}
  \hline
  \textbf{Year} & \textbf{Qu} & \textbf{Topic} & \textbf{What it tests} & \textbf{Ans}\\ \hline
  \endhead
    2019 & 2 & T1 & Two stars' line spectra: relative motion (shift) and composition & B\\ \hline
    2019 & 6 & T3 & Graph of peak wavelength against black-body temperature (inverse) & A\\ \hline
    2019 & 10 & T2 & Malus' law with polarisers at 30$^\circ$: intensity ratio and model of light & A\\ \hline
    2019 & 13 & T3/\allowbreak M8 & Photons emitted per second by a laser of given power and wavelength & B\\ \hline
    2019 & 15 & T2 & Double slit: correct slit spacing and angle in $d\sin\theta = m\lambda$ & C\\ \hline
    2020 & 1 & T2 & Who proposed the wavefront model that explains refraction & D\\ \hline
    2020 & 3 & T1 & What Maxwell's prediction of the speed of EM waves was based on & C\\ \hline
    2020 & 13 & T3 & $K_\mathrm{max}$-frequency lines for two metals: compare maximum photoelectron speeds & A\\ \hline
    2020 & 17 & T4 & Light pulse from a $0.95c$ train: speed, distance and time in both frames & C\\ \hline
    2020 & 18 & T1 & Romer's Io-eclipse timing: Earth position giving the longest observed period & B\\ \hline
    2021 & 4 & T4 & Ship at $0.8c$: which quantity both observers agree on (speed of light) & B\\ \hline
    2021 & 5 & T1/\allowbreak M8 & Identify the source (star) and line type (absorption) from a spectrum & A\\ \hline
    2021 & 8 & T2 & Bright spot at the centre of a disc's shadow: wave effects responsible & D\\ \hline
    2021 & 11 & T3 & Wien's law: peak wavelength emitted by a human body at 310 K & A\\ \hline
    2021 & 15 & T2 & Two polarisers, first one rotated: which transmitted intensities change & C\\ \hline
    2021 & 16 & T4 & Sun's mass loss per minute from its power output ($E = mc^2$) & B\\ \hline
    2021 & 20 & T3/\allowbreak M6 & Photoelectrons circling in a magnetic field: radius when photon energy doubles & D\\ \hline
    2022 & 2 & T1 & Which property of a star produces absorption lines & D\\ \hline
    2022 & 7 & T3 & Photon frequency from its energy ($E = hf$) & C\\ \hline
    2022 & 9 & T3 & Black-body temperature from its peak wavelength (Wien) & C\\ \hline
    2022 & 14 & T3 & $K_\mathrm{max}$-frequency graph: change that gives a parallel, shifted line & B\\ \hline
    2022 & 17 & T2 & Rotating polariser before crossed polarising lenses: compare intensities & C\\ \hline
    2022 & 20 & T4 & Wheel rolling at $0.4c$: length contraction at the bottom, centre and top & B\\ \hline
    2023 & 3 & T2 & Expected screen pattern from a double-slit set-up & C\\ \hline
    2023 & 6 & T1 & Condition for an electron to radiate EM waves (acceleration) & D\\ \hline
    2023 & 9 & T3 & How a black-body curve changes when the temperature falls from 4500 K to 4000 K & A\\ \hline
    2023 & 15 & T3 & Photoelectric observation that supports Einstein's photon model & A\\ \hline
    2023 & 19 & T4/\allowbreak M6 & Charge spacing in current-carrying wires seen from a moving charge's frame & B\\ \hline
    2024 & 2 & T2 & Evidence that supports Huygens' wave model of light & B\\ \hline
    2024 & 6 & T3 & Meaning of the work-function term in $K_\mathrm{max} = hf - \phi$ & C\\ \hline
    2024 & 12 & T4 & Curve of measured rod length against speed (length contraction) & B\\ \hline
    2024 & 16 & T3 & $K_\mathrm{max}$ against photon energy for four metals: which emit at a given frequency & A\\ \hline
    2024 & 20 & T4 & Clocks on the ground, in a 12-h orbit and geostationary: tick rates (SR only) & B\\ \hline
    2025 & 1 & T1 & Maxwell's contribution to the understanding of light & B\\ \hline
    2025 & 9 & T1 & Stellar property (density) from absorption-line widths & A\\ \hline
    2025 & 12 & T2 & Graph of intensity against $\cos^2\theta$ that represents Malus' law & A\\ \hline
    2025 & 19 & T4 & Mass of a sealed system after an energy-releasing reaction & B\\ \hline
\end{longtable}}

Summary: 37 questions, 4 to 7 in each paper (T3: 12, T2: 9, T1: 8, T4: 8).
The quantum model (black-body radiation, Wien's law and the photoelectric effect, especially $K_\mathrm{max}$ graphs) is tested most, then the wave model (double slits and Malus' law). Special relativity appears in every paper from 2020, and stellar spectra or Maxwell's work in most years.
\end{document}
